\makeabstract
{
Investigating Electronic and Steric Modulation in Group V Supported Titanium Alkoxides for Controlled Ring-Opening Polymerization
}
{
Trey Sheppard (1), Greta Stechschulte (2), Christopher B. Durr, PhD (1,2,3)
}
{
\textsuperscript{1}Amherst College, Amherst, MA, USA\\
\textsuperscript{2}The Ohio State University, Columbus, OH, USA
}
{
To determine how electronic donor strength and steric bulk of (Group IV) donor scaffolds on Ti(IV) alkoxide initiators govern initiation/propagation kinetics and polymerization control in the ring-opening polymerization of cyclic esters, so we can establish design rules for living, predictable ROP that yields biodegradable polymers.
}
{
Ligands were assembled from commercially available precursors using literature-adapted routes and prior in-lab protocols. Structures (and purity) were confirmed by NMR (¹H; others as applicable) and by SCXRD when suitable crystals were obtained. Air-stable ligand steps were performed in a fume hood; air-sensitive Ti (IV) alkoxide complexes were synthesized and handled exclusively in an anhydrous, O₂-free glovebox. Ring-opening polymerizations of cyclic esters were run neat (melt conditions) by combining monomer and catalyst at 140 °C inside the glovebox.

}
{

We crystallized several titanium complexes and found a binding geometry that differs from prior reports. Their structures were confirmed by SCXRD and supported by solution NMR. In melt ROP, a representative Ti catalyst achieved 97\% monomer conversion in 1h and produced a polymer with a DP ≈ of 75. We will use size-exclusion chromatography to measure molecular weights and dispersity, assessing polymer quality. Overall, these results indicate a promising Ti platform, allowing us to compare the effects of different substituents on structure and performance.

}
{

We examined the kinetics and performance of titanium alkoxide catalysts for ring-opening polymerization of cyclic esters. By varying ligand functional groups, we directly compared how changes in donor electronics and steric bulk affect catalyst activity and control. The modular ligand set provides many new derivatives to explore next. Overall, the data support continued development of Ti-based systems to produce biodegradable polymers with targeted properties and predictable control.
}

\newpage
